\documentclass[11pt,a4paper]{article}
\usepackage[T1]{fontenc}
\usepackage{lmodern,amsmath,amssymb,amsthm,microtype}
\usepackage[margin=27mm]{geometry}
\usepackage[colorlinks=true,urlcolor=blue]{hyperref}
\setlength{\emergencystretch}{2em}
\newtheorem{proposition}{Proposition}
\newcommand{\Q}{\mathbb Q}
\title{Scope of the Achenjang--Woo prime and twist theorems}
\author{}\date{17 November 2026}
\begin{document}\maketitle

\section{The cited statements}

The source is the complete 74-page manuscript \cite{AW}, version 1,
submitted 11 November 2026, with manuscript date 13 November 2026.
Theorem A, on page 3, treats a fixed system of pairwise independent
binary affine linear forms together with one positive definite binary
quadratic norm form from an imaginary quadratic field. It gives the
prime-value asymptotic, and simultaneous prime values when the singular
series is positive. The norm form may have the separate affine changes
$N_{K/\Q}(T_1(x)+T_2(y)\alpha)$ displayed in that theorem.
It does not state a simultaneous-prime theorem for several independent
irreducible binary quadratic forms.

Theorem 10.1, on page 39, gives infinitely many rank-exactly-two
quadratic twists of
\[
 E_{\mu,\lambda}:y^2=x(x-\mu)(x-\lambda)
\]
when $(\mu,\lambda)$ satisfies Definition 10.4, page 40. The
conditions include $\mu$ odd and squarefree, $\lambda=\mu m^2$,
$m\equiv2\pmod8$, and the following two essential auxiliary-prime
requirements:
\begin{itemize}
\item two distinct primes $p_1,p_2>14$ divide $m^2-1$ to odd order,
      with the additional Hilbert-symbol and residue requirements
      specified in (A3);
\item two distinct odd primes $\ell_0,\ell_1$ divide $m$ and do not
      divide $\mu$, with the additional splitting requirement in (A4).
\end{itemize}
The coprimality and parity conditions (A5)--(A6) are also required.
The theorem is not a result for all curves with three rational
two-torsion points. Its statement does not prescribe arbitrary
abscissa square classes or an external polynomial family of twist
parameters.

\section{The original pole family}

\begin{proposition}
Theorem 10.1 does not apply directly to any quadratic twist of the
original pole curve
\[
 E_n:y^2=x(x-n)(x+3n),\qquad n\ne0,
\]
in any ordering of its three rational roots.
\end{proposition}
\begin{proof}
The substitution $X=x-n$ gives
\[
 y^2=X(X+n)(X+4n),
\]
so the pair of nonzero roots can be written $(\mu,\lambda)=(-n,-4n)$.
The ratio is $4$, giving $m=2$. But $m^2-1=3$ has no prime factor
greater than fourteen, and $m$ has no odd prime factor. Thus both
(A3) and (A4) fail, independently of $n$.

The possible ratios from all root orderings are
\[
 -3,\quad -\tfrac13,\quad 4,\quad \tfrac14,
 \quad\tfrac34,\quad\tfrac43.
\]
The only square ratios are $4$ and $1/4$; reordering to make the
square root a positive integer gives $m=2$ again. A quadratic twist
does not change these ratios. Hence it cannot repair these failures.
\end{proof}

The same source also does not apply directly to the intrinsic
congruent-number family $y^2=x(x-n)(x+n)$: its root ratios are
$-1,2,1/2$, none of which is a rational square.

\section{The generic two-point family and the missing incidence}

Proposition 10.2 of \cite{AW} recalls the Rubin--Silverberg family.
For distinct nonzero $\mu,\lambda$ with $\mu+\lambda\ne0$, put
\[
 \alpha=\frac{\mu\lambda}{\mu+\lambda},\qquad
 G(u)=\frac{(u^2+\alpha^2)(u^2+\alpha^2-\alpha\mu)
                    (u^2+\alpha^2-\alpha\lambda)}{\alpha^3}.
\]
The twist $G(u)y^2=x(x-\mu)(x-\lambda)$ has the two independent
generic sections
\[
 \left(\frac{u^2+\alpha^2}{\alpha},1\right),\qquad
 \left(\frac{\alpha u^2+\alpha^3}{u^2},\frac{\alpha^3}{u^3}\right).
\]
This construction remains valid for $(\mu,\lambda)=(-1,-4)$,
although the exact-rank theorem does not apply there.
Putting $u=2v/5$ and translating back to the original root ordering
gives
\[
 G=-\frac{(v^2+4)(v^2-1)(v^2-16)}{125}
\]
and points on $E_G$ with
\begin{align*}
 x_1&=G\frac{1-v^2}{5}, & y_1&=G^2,\\
 x_2&=G\frac{v^2-16}{5v^2}, & y_2&=-\frac{8G^2}{v^3}.
\end{align*}
Their abscissa classes, away from the indicated zeros and poles, are
\[
 [x_1]=[(v^2+4)(v^2-16)],\qquad
 [x_2]=[-(v^2+4)(v^2-1)]
 \quad\text{in }\Q^\times/\Q^{\times2}.
\]
These formulas give individual rank-at-least-two twist families.
They do not prescribe arbitrary classes $D_i$ on the native curves
\[
 E_{N_i(\tau)},\qquad
 N_i(\tau)=D_iC_i(\tau^2-D_i)(\tau^2+3D_i),
\]
and do not furnish a common rational $\tau$.
Matching a specialization would require, at a minimum,
$G(v_i)/N_i(\tau)\in\Q^{\times2}$ and the required nonzero point
class. Neither condition follows from generic rank two or from
the prime theorem. The relation must also use the actual $D_i,C_i$
and actual detector traces.

\section{Scope for the single-square prime proposal}

In the proposal $t=A+z^2$, a common square denominator for $z$ makes
the native contact factors proportional to
\[
 m^2+(A-\theta_j)n^2.
\]
For the source's large positive $A$ and distinct fixed real contact
roots $\theta_j$, these are distinct positive definite irreducible
binary quadratic forms. Theorem A of \cite{AW} permits one such norm
form alongside linear forms; it does not permit this whole collection.
Thus it does not fill the restricted-prime-pattern gap left by Sun's
free-translation theorem.

The conic-bundle consequence, Theorem B and its generalization
Theorem 9.2, concerns sufficiency of the Brauer--Manin obstruction
for weak approximation in the displayed family. It cannot be used
without verifying that a proposed incidence variety belongs to that
family and has an appropriate adelic point orthogonal to its Brauer
group. Neither verification has been made for the detector incidence.

\begin{thebibliography}{9}\raggedright
\bibitem{AW}N.~Achenjang and K.~Woo,
\emph{Prime values of linear forms and a quadratic form: one proof
and two arithmetic applications}, arXiv:2611.16062v1 (2026),
\url{https://arxiv.org/abs/2611.16062}.
The full author manuscript was inspected through
\url{https://arxiv.cnoctave.top/pdf/2611.16062};
Theorem A, Theorem 10.1, Proposition 10.2 and Definition 10.4.
\end{thebibliography}
\end{document}
